\subsection*{11.4 Application: Data Compression}

Let A ={ a1,a 2,...,a m} denote an alphabet of size m. Suppose we generate

a string of length n, with each character drawn independently according to a

probability distribution P(ai) = pi , where pi 's are nonnegative real numbers

summing to 1. There are totally mn combinations. We may associate each string

with an integer between 1 and mnU s i n g log2(mn) bits, we can store and recover

the random string with no decoding error.

However, if we allow an arbitrarily small error in decoding, we can significantly

reduce the number of bits. It turns out there exists a set S\subset A n with probability

Pr(S)\geq 1 - \epsilon and \vertS\vert\ll mn. We only encode the strings in the set S The strings

not in S are not encoded, or encoded arbitrarily. The decoding error is thus less than

or equal to \epsilon. We will show that \epsiloncan be made arbitrarily small as the block length

n increases, with the cardinality of the set S roughly equal to .2nH(P) , where H(P)

denotes the entropy of the distribution P ,

H(P) \coloneqq-

m\sum

i=1

pi log2 pi.

To realize this idea, we fix a parameter \delta> 0 and define a set S\delta,n of strings by

S\delta,n\coloneqq{ (x1,x 2,...,x n)\inA n :2 -n(H(P) +\delta) \leq

n

k=1

p(xk)\leq 2 -n(H(P) -\delta)}.

We claim that for any fixed \delta> 0,t h e s e t S\delta,n has probability arbitrarily close to 1

asn\to\infty We denote the random string of length n by. X, and for k= 1 ,2,...,n ,

let Xk be the k -th letter in X. Define a new random variable Uk =- log2 P(Xk),

where P(Xk) is equal to pi if Xk =i For simplicity, we may assume pi >0 for

all i , so that log2 pi is well-defined. The expectation of Uk is precisely equal to the

entropy H(P) of distribution P By the weak law of large numbers, for any fixed

\delta> 0,we have

limn\to\infty Pr

( \vert\vert\vert 1

n

n\sum

k=1

Uk -H(P)

\vert\vert\vert \leq\delta

)

=1 .

But

\vert\vert\vert 1

n

n\sum

k=1

Uk -H(P)

\vert\vert\vert \leq\delta 2 n(H(P) -\delta) \leq2 U1+U2+\cdot\cdot\cdot+Un \leq2 n(H(P) +\delta)

2 -n(H(P) +\delta) \leqP(X 1)P(X 2)\cdot\cdot\cdot P(Xn)\leq 2 -n(H(P) -\delta).

The random string X belongs toS\delta,n if and only if .

\vert\vert\vert 1

n

\sumn

k=1 Uk -H(P)

\vert\vert\vert \leq\deltaT h i s

shows that the probability of the set of strings S\delta,n approaches 1 asn\to\infty , for any

fixed \delta> 0. To estimate the cardinality of S\delta,n, we note that all random strings in

S\delta,n have probability at least .2-n(H(P) +\delta)There are at most .2n(H(P) +\delta) strings in

S\delta,n. We can represent them using at most n(H(P) + \delta) bits.

To design the data compression scheme with decoded error no more than \epsilon,w e

first fix a small \delta> 0 and find a sufficiently large block length n such that Pr(X\in

S\delta,n)> 1 - \epsilon. Because we only encode the strings in S\delta,n,it takes H(P) + \deltabits

per symbol. The saving is particularly significant if H(P) \ll log2(m).

The result demonstrated in this example is the forward part of Shannon's theorem

on data compression. This is one of the first major results in information theory [ 3].

We demonstrate the idea of this data compression with finite block length n.

Suppose that alphabet is A ={ a,b,c }The data source generates symbols a, b,

and c independently, according to the distribution P(a) = 2/3, P(b) = 2/9 and

P(c) = 1/9. We group the random symbols into groups of size n = 11 and use

B bits to encode each group. Only the .2B - 1 strings of length n with the largest

probabilities are encoded. If the random string is not among these .2B - 1 selected

ones, we use a special number, such as . 2B , to signify that a decoding error has

occurred.

The probability of error is calculated by the following Python program.

Python Program for Demonstrating Data Compression

# Compute the success decoding probability of

# data compression with fixed block length

from numpy import prod

from itertools import product

n = 11 # block length

prob_dist= {'a':2/3, 'b':2/9, 'c':1/9}

B=1 6

A = [x for x in prob_dist] # list of letters

D={}

for t in product(A,repeat = n):

s = ''join(t)

D[s] = prod([prob_dist[x] for x in t])

# sum the largest 2 **B-1 probabilities

list_of_prob = list(Dvalues())

list_of_probsort(reverse = True)

prob_success = sum(list_of_prob[0:(2 **B-1)])

print(f"{B} bits per block")

print(f"P(successful decoding) = {prob_success}")

The table below shows the probabilities of successful decoding for a different

number of bits per block. The probabilities are obtained by running the Python

program forB = 8, 9,..., 16.

No. of bits per block 8 9 10 11 12 13 14 15 16

P (successful decoding) 0.245 0.327 0.428 0.533 0.650 0.765 0.864 0.939 0.983

We observe that the value of .11 \cdot H(P) = 11(1.22) 13.475, where H(P)

1.22 is entropy of the probability distribution P The table shows that there is a

significant increase in the probability of successful decoding when the number of

bits per block is slightly above the threshold value, ie., for B= 14 or higher.

In general, as the block length n approaches infinity, there is an abrupt increase

in success probability when the number of bits per block is .1.22n.
